\documentclass[11pt]{article}
\usepackage{reusv}
\usepackage{listings}
\usepackage{enumitem}

\lstset{basicstyle=\ttfamily\small, frame=single, breaklines=true}
\setborderthickness{0.5pt}
\setbordermargin{1cm}
\setbordercolor{black}

\slimtitle{C1. NLP 정식화 및 Interior-Point 솔버}

\begin{document}
\drawborder
\thispagestyle{firstpage}

\lightertitle{C1. NLP 정식화 및 Interior-Point 솔버}
\def\lighterdate{February 12, 2026}
\lighterauthor{Suwon Lee}
\lighteraddress{}{Future Mobility Control Lab, Department of Future Mobility, Kookmin University, Seoul, Republic of Korea}
\lighteremail{suwon@kookmin.ac.kr}

\begin{abstract}
연속 추력 궤도전이 문제는 본래 무한 차원 함수 공간에서 정의된 최적 제어 문제(OCP)이다.
본 문서는 OCP에서 유한 차원 비선형 프로그래밍(NLP)으로의 변환 과정,
interior-point 솔버(IPOPT)의 수학적 원리(장벽 함수, KKT 조건, 프라이멀-듀얼 Newton 시스템, 적응적 장벽 파라미터 전략),
그리고 Two-Pass 전략과 수렴 실패 복구 메커니즘을 다룬다.
구현은 CasADi의 Opti stack을 통해 NLP를 심볼릭으로 구성하고 IPOPT 솔버에 전달하는 구조이다.
\end{abstract}

\vspace{0.3cm}
\noindent\textit{대응 코드:} \texttt{src/orbit\_transfer/optimizer/solver.py}, \texttt{src/orbit\_transfer/optimizer/two\_pass.py}

%% ============================================================================
\section{목적}
%% ============================================================================

연속 추력 궤도전이 문제는 본래 무한 차원 함수 공간에서 정의된 최적 제어 문제(Optimal Control Problem, OCP)이다. 상태 $\mathbf{x}(t) \in \mathbb{R}^6$과 제어 $\mathbf{u}(t) \in \mathbb{R}^3$은 시간의 연속 함수이며, 동역학은 상미분방정식(ODE) $\dot{\mathbf{x}} = \mathbf{f}(\mathbf{x}, \mathbf{u})$로 기술된다. 이러한 OCP를 수치적으로 풀기 위해서는 유한 차원의 비선형 프로그래밍(Nonlinear Programming, NLP) 문제로 변환하는 이산화(transcription) 과정이 필요하다.

본 문서는 다음 세 가지 주제를 다룬다.

\begin{enumerate}[leftmargin=*]
  \item \textbf{OCP에서 NLP로의 변환}: Direct collocation에 의해 연속 시간 문제를 유한 차원 NLP 표준형으로 매핑하는 과정
  \item \textbf{Interior-point 솔버(IPOPT)의 수학적 원리}: 장벽 함수, KKT 조건, 프라이멀-듀얼 Newton 시스템, 적응적 장벽 파라미터 전략
  \item \textbf{Two-Pass 전략과 수렴 실패 복구}: 허용 오차 연속체(continuation), 다중 시작점 전략, 실행가능성 복원
\end{enumerate}

구현은 CasADi의 Opti stack을 통해 NLP를 심볼릭으로 구성하고, IPOPT 솔버에 전달하는 구조이다. \texttt{solver.py}는 IPOPT 옵션 관리를, \texttt{two\_pass.py}는 Two-Pass 최적화 파이프라인의 전체 흐름과 수렴 실패 복구 로직을 담당한다.


%% ============================================================================
\section{수학적 배경}
%% ============================================================================

\subsection{최적 제어 문제의 정식화}

시간 구간 $[0, T]$에서 다음 볼차(Bolza) 형태의 최적 제어 문제를 고려한다.
%
\begin{equation}\label{eq:ocp-cost}
\min_{\mathbf{u}(t),\, \nu_0,\, \nu_f} \quad J = \int_0^T \|\mathbf{u}(t)\|^2 \, dt
\end{equation}
%
subject to:
%
\begin{equation}\label{eq:ocp-dynamics}
\dot{\mathbf{x}}(t) = \mathbf{f}(\mathbf{x}(t),\, \mathbf{u}(t)), \quad t \in [0, T]
\end{equation}
%
\begin{equation}\label{eq:ocp-bc}
\mathbf{x}(0) = \mathbf{g}_0(\nu_0), \quad \mathbf{x}(T) = \mathbf{g}_f(\nu_f)
\end{equation}
%
\begin{equation}\label{eq:ocp-ineq}
\|\mathbf{u}(t)\| \leq u_{\max}, \quad \|\mathbf{r}(t)\| \geq R_E + h_{\min}
\end{equation}

여기서 $\mathbf{x} = [\mathbf{r}^\top,\, \mathbf{v}^\top]^\top \in \mathbb{R}^6$은 ECI 좌표계의 위치-속도 상태벡터, $\mathbf{u} \in \mathbb{R}^3$은 추력 가속도 벡터이다. $\mathbf{f}$는 J2 섭동을 포함한 운동방정식이며, $\mathbf{g}_0(\nu_0)$과 $\mathbf{g}_f(\nu_f)$는 궤도요소에서 ECI 상태벡터로의 변환 함수이다. $\nu_0$와 $\nu_f$는 출발/도착 궤도상의 위치를 나타내는 자유변수로, 솔버가 최적의 출발-도착 위치를 자동으로 결정한다.

\subsection{NLP 표준형}

일반적인 NLP 문제는 다음과 같이 정의된다~\cite{Nocedal2006}.
%
\begin{equation}\label{eq:nlp-obj}
\min_{\mathbf{z} \in \mathbb{R}^n} \quad f(\mathbf{z})
\end{equation}
%
\[
\text{s.t.} \quad \mathbf{c}(\mathbf{z}) = \mathbf{0} \quad (\mathbf{c} : \mathbb{R}^n \to \mathbb{R}^{m_e})
\]
%
\[
\quad\quad\;\; \mathbf{d}(\mathbf{z}) \geq \mathbf{0} \quad (\mathbf{d} : \mathbb{R}^n \to \mathbb{R}^{m_i})
\]

여기서 $\mathbf{z}$는 결정변수 벡터, $f$는 목적함수, $\mathbf{c}$는 등식 제약, $\mathbf{d}$는 부등식 제약이다. $n$은 결정변수의 수, $m_e$는 등식 제약의 수, $m_i$는 부등식 제약의 수이다.

\subsection{OCP에서 NLP로의 변환}

Direct collocation 이산화를 적용하면, 연속 시간 OCP가 유한 차원 NLP로 변환된다. 결정변수 벡터는 다음과 같이 구성된다.
%
\begin{equation}\label{eq:decision-vec}
\mathbf{z} = \begin{bmatrix} \text{vec}(\mathbf{X}) \\ \text{vec}(\mathbf{U}) \\ \nu_0 \\ \nu_f \end{bmatrix} \in \mathbb{R}^n
\end{equation}

여기서 $\mathbf{X} \in \mathbb{R}^{6 \times N}$은 $N$개 배치점에서의 상태, $\mathbf{U} \in \mathbb{R}^{3 \times N}$은 제어 배열이다. $\text{vec}(\cdot)$은 행렬을 열 벡터로 펼치는 연산이다.

\textbf{목적함수.} 연속 시간 비용함수 $J = \int_0^T \|\mathbf{u}\|^2\, dt$는 수치 적분 규칙에 의해 이산 비용함수 $f(\mathbf{z})$로 근사된다. Hermite-Simpson collocation에서는 Simpson 구적법을, Multi-Phase LGL collocation에서는 LGL 가중치 적분을 사용한다.

\textbf{등식 제약 $\mathbf{c}(\mathbf{z}) = \mathbf{0}$.} 등식 제약은 두 부류로 구성된다.

\begin{enumerate}[leftmargin=*]
  \item \textbf{Collocation defect 제약}: 각 부분 구간에서 이산화된 동역학 방정식의 잔차가 0이 되어야 한다는 조건이다. Hermite-Simpson의 경우 Simpson 연속성 조건과 Hermite 중간점 조건이, LGL의 경우 미분 행렬에 의한 collocation 조건과 phase간 linkage 조건이 이에 해당한다.

  \item \textbf{경계조건}: 출발/도착 궤도의 ECI 상태벡터가 궤도요소로부터 결정된 값과 일치해야 한다.
\end{enumerate}
%
\begin{equation}\label{eq:bc-constraint}
\mathbf{X}_{:,0} - \mathbf{g}_0(\nu_0) = \mathbf{0}, \quad \mathbf{X}_{:,-1} - \mathbf{g}_f(\nu_f) = \mathbf{0}
\end{equation}

\textbf{부등식 제약 $\mathbf{d}(\mathbf{z}) \geq \mathbf{0}$.} 부등식 제약은 각 배치점에서 부과된다.
%
\begin{equation}\label{eq:thrust-bound}
u_{\max}^2 - \mathbf{u}_j^\top \mathbf{u}_j \geq 0, \quad j = 0, 1, \ldots, N-1 \quad \text{(추력 상한)}
\end{equation}
%
\begin{equation}\label{eq:altitude-bound}
\mathbf{r}_j^\top \mathbf{r}_j - (R_E + h_{\min})^2 \geq 0, \quad j = 0, 1, \ldots, N-1 \quad \text{(최소 고도)}
\end{equation}

\subsection{NLP의 규모}

Pass~1(Hermite-Simpson, $M = 30$)에서의 NLP 규모는 다음과 같다.

\begin{table}[h]
\centering
\begin{tabular}{@{}ll@{}}
\toprule
구성요소 & 크기 \\
\midrule
결정변수 $n$ & $6 \times 61 + 3 \times 61 + 2 = 551$ \\
등식 제약 $m_e$ & $6 \times 30 + 6 \times 30 + 6 + 6 = 372$ \\
부등식 제약 $m_i$ & $61 + 61 = 122$ \\
제약 합계 & 494 \\
\bottomrule
\end{tabular}
\end{table}

Pass~2(Multi-Phase LGL)에서의 NLP 규모는 phase 구조에 따라 달라진다. 예를 들어 bimodal 프로파일에서 3개 phase(피크 15노드 + coasting 8노드 + 피크 15노드)를 사용하는 경우, 결정변수는 $6 \times 38 + 3 \times 38 + 2 = 344$개이다. Phase 경계에서의 linkage 제약이 추가되는 반면, 전체 노드 수가 줄어들어 collocation 제약 수는 감소한다.

\subsection{KKT 최적성 조건}

NLP 표준형의 라그랑지안(Lagrangian)을 다음과 같이 정의한다.
%
\begin{equation}\label{eq:lagrangian}
\mathcal{L}(\mathbf{z}, \boldsymbol{\lambda}, \boldsymbol{\mu}) = f(\mathbf{z}) - \boldsymbol{\lambda}^\top \mathbf{c}(\mathbf{z}) - \boldsymbol{\mu}^\top \mathbf{d}(\mathbf{z})
\end{equation}

여기서 $\boldsymbol{\lambda} \in \mathbb{R}^{m_e}$는 등식 제약의 라그랑주 승수, $\boldsymbol{\mu} \in \mathbb{R}^{m_i}$는 부등식 제약의 라그랑주 승수이다.

Karush-Kuhn-Tucker(KKT) 조건은 정칙성 가정(constraint qualification) 하에서 국소 최적해의 필요조건이다~\cite{Nocedal2006}.

\textbf{(1) 정상성 조건 (Stationarity):}
%
\begin{equation}\label{eq:kkt-stationarity}
\nabla_{\mathbf{z}} \mathcal{L} = \nabla f(\mathbf{z}) - \nabla \mathbf{c}(\mathbf{z})^\top \boldsymbol{\lambda} - \nabla \mathbf{d}(\mathbf{z})^\top \boldsymbol{\mu} = \mathbf{0}
\end{equation}

\textbf{(2) 프라이멀 실행가능성 (Primal feasibility):}
%
\begin{equation}\label{eq:kkt-primal}
\mathbf{c}(\mathbf{z}) = \mathbf{0}, \quad \mathbf{d}(\mathbf{z}) \geq \mathbf{0}
\end{equation}

\textbf{(3) 듀얼 실행가능성 (Dual feasibility):}
%
\begin{equation}\label{eq:kkt-dual}
\boldsymbol{\mu} \geq \mathbf{0}
\end{equation}

\textbf{(4) 상보성 조건 (Complementarity):}
%
\begin{equation}\label{eq:kkt-complementarity}
\mu_j \, d_j(\mathbf{z}) = 0, \quad j = 1, \ldots, m_i
\end{equation}

상보성 조건은 각 부등식 제약에 대해 ``승수가 0이거나 제약이 등호로 활성(active)이다''라는 것을 의미한다. 이 조건들을 모두 벡터로 모으면 KKT 시스템을 얻는다. Interior-point 방법은 이 KKT 시스템의 완화된(perturbed) 버전을 Newton 방법으로 풀어나가는 전략이다.

\subsection{Interior-Point (장벽) 방법}

Interior-point 방법은 부등식 제약 $\mathbf{d}(\mathbf{z}) \geq \mathbf{0}$을 로그 장벽 함수(logarithmic barrier function)로 대체하여, 부등식 제약이 없는 등식 제약 문제의 수열로 변환하는 방법이다~\cite{Fiacco1968, Wachter2006}.

\textbf{장벽 부문제(Barrier subproblem).} 장벽 파라미터 $\mu > 0$에 대해 다음 문제를 정의한다.
%
\begin{equation}\label{eq:barrier-obj}
\min_{\mathbf{z}} \quad \varphi_\mu(\mathbf{z}) = f(\mathbf{z}) - \mu \sum_{j=1}^{m_i} \ln\bigl(d_j(\mathbf{z})\bigr)
\end{equation}
%
\[
\text{s.t.} \quad \mathbf{c}(\mathbf{z}) = \mathbf{0}
\]

로그 장벽 항 $-\mu \ln(d_j)$는 $d_j \to 0^+$일 때 $+\infty$로 발산하여, 해가 실행가능 영역의 내부(interior)에 머물도록 강제한다. $\mu \to 0^+$으로 감소시키면 장벽 부문제의 해가 원래 NLP의 해에 수렴한다.

\textbf{완화된 KKT 시스템.} 장벽 부문제의 KKT 조건은 원래 KKT 조건에서 상보성 조건만 변경된다.
%
\[
\nabla f(\mathbf{z}) - \nabla \mathbf{c}(\mathbf{z})^\top \boldsymbol{\lambda} - \nabla \mathbf{d}(\mathbf{z})^\top \boldsymbol{\mu} = \mathbf{0}
\]
%
\[
\mathbf{c}(\mathbf{z}) = \mathbf{0}
\]
%
\begin{equation}\label{eq:perturbed-complementarity}
\mu_j \, d_j(\mathbf{z}) = \mu, \quad j = 1, \ldots, m_i \quad \text{(완화된 상보성)}
\end{equation}
%
\[
\boldsymbol{\mu} > \mathbf{0}, \quad \mathbf{d}(\mathbf{z}) > \mathbf{0}
\]

원래 상보성 조건 $\mu_j d_j = 0$이 $\mu_j d_j = \mu$로 완화되었다. $\mu \to 0$이면 원래의 상보성 조건이 복원된다.

\subsection{프라이멀-듀얼 Newton 시스템}

완화된 KKT 시스템을 Newton 방법으로 선형화하면, 각 반복에서 다음 선형 시스템을 풀어야 한다. 슬랙 변수 $\mathbf{s} = \mathbf{d}(\mathbf{z})$를 도입하면, 프라이멀-듀얼 Newton 시스템은 다음과 같다~\cite{Wachter2006}.
%
\begin{equation}\label{eq:newton-full}
\begin{bmatrix}
\mathbf{W} & -\nabla \mathbf{c}^\top & -\nabla \mathbf{d}^\top & \mathbf{0} \\
\nabla \mathbf{c} & \mathbf{0} & \mathbf{0} & \mathbf{0} \\
\nabla \mathbf{d} & \mathbf{0} & \mathbf{0} & -\mathbf{I} \\
\mathbf{0} & \mathbf{0} & \mathbf{S} & \mathbf{M}
\end{bmatrix}
\begin{bmatrix}
\Delta \mathbf{z} \\
\Delta \boldsymbol{\lambda} \\
\Delta \boldsymbol{\mu} \\
\Delta \mathbf{s}
\end{bmatrix}
= -
\begin{bmatrix}
\nabla_\mathbf{z} \mathcal{L} \\
\mathbf{c} \\
\mathbf{d} - \mathbf{s} \\
\mathbf{S} \mathbf{M} \mathbf{e} - \mu \mathbf{e}
\end{bmatrix}
\end{equation}

여기서

\begin{itemize}[leftmargin=*]
  \item $\mathbf{W} = \nabla^2_{zz} \mathcal{L}$: 라그랑지안의 헤시안
  \item $\mathbf{S} = \text{diag}(s_1, \ldots, s_{m_i})$, $\mathbf{M} = \text{diag}(\mu_1, \ldots, \mu_{m_i})$
  \item $\mathbf{e} = [1, \ldots, 1]^\top$
\end{itemize}

$\Delta \mathbf{s}$를 소거하면, 축소된(condensed) 시스템을 얻는다.
%
\begin{equation}\label{eq:newton-condensed}
\begin{bmatrix}
\mathbf{W} + \nabla \mathbf{d}^\top \mathbf{S}^{-1} \mathbf{M} \nabla \mathbf{d} & -\nabla \mathbf{c}^\top \\
\nabla \mathbf{c} & \mathbf{0}
\end{bmatrix}
\begin{bmatrix}
\Delta \mathbf{z} \\
\Delta \boldsymbol{\lambda}
\end{bmatrix}
= -
\begin{bmatrix}
\nabla_\mathbf{z} \mathcal{L} + \nabla \mathbf{d}^\top (\boldsymbol{\mu} - \mu \mathbf{S}^{-1} \mathbf{e}) \\
\mathbf{c}
\end{bmatrix}
\end{equation}

이 축소 시스템은 대칭이지만 부정치(indefinite)인 행렬을 갖는다. IPOPT는 이 시스템을 직접 솔버(예: MUMPS)로 인수분해하여 풀며, 관성 보정(inertia correction)을 통해 수치적 안정성을 유지한다.

\subsection{IPOPT의 적응적 장벽 파라미터 전략}

장벽 파라미터 $\mu$의 감소 전략은 interior-point 방법의 수렴 속도를 결정짓는 핵심 요소이다. IPOPT는 두 가지 전략을 제공한다~\cite{Wachter2006}.

\textbf{(1) 단조 감소(monotone) 전략.} 고전적인 접근법으로, 장벽 부문제가 충분히 수렴하면 $\mu$를 고정 비율로 감소시킨다.
%
\begin{equation}\label{eq:mu-monotone}
\mu_{k+1} = \kappa_\mu \, \mu_k, \quad 0 < \kappa_\mu < 1
\end{equation}

\textbf{(2) 적응적(adaptive) 전략.} 현재 상보성 잔차에 기반하여 $\mu$를 적응적으로 설정한다.
%
\begin{equation}\label{eq:mu-adaptive}
\mu_{k+1} = \max\left(\frac{\epsilon_{\text{tol}}}{10},\; \min\left(\kappa_\mu \, \mu_k,\; \mu_k^{\theta_\mu}\right)\right)
\end{equation}

여기서 $\theta_\mu > 1$은 초선형 수렴을 유도하는 지수이다. 적응적 전략은 $\mu$가 빠르게 감소하지만 실행가능 영역을 벗어나지 않는 균형점을 추구한다.

본 구현에서는 \texttt{mu\_strategy = "adaptive"}를 사용한다. 적응적 전략은 궤적 최적화와 같이 제약조건의 비선형성이 강한 문제에서 단조 전략보다 안정적인 수렴을 보이는 것으로 알려져 있다~\cite{Betts2010}.

\subsection{MUMPS 직접 솔버}

각 Newton 반복에서 축소된 KKT 시스템(2.7절)의 계수 행렬은 대칭 부정치 희소 행렬(sparse symmetric indefinite matrix)이다. MUMPS(MUltifrontal Massively Parallel Sparse direct Solver)는 이러한 행렬에 대해 Bunch-Kaufman 피벗팅 기반의 $\mathbf{LDL}^\top$ 분해를 수행한다~\cite{Amestoy2001}.
%
\begin{equation}\label{eq:ldlt}
\mathbf{K} = \mathbf{P} \mathbf{L} \mathbf{D} \mathbf{L}^\top \mathbf{P}^\top
\end{equation}

여기서 $\mathbf{P}$는 순열 행렬, $\mathbf{L}$은 하삼각 행렬, $\mathbf{D}$는 $1 \times 1$ 및 $2 \times 2$ 블록 대각 행렬이다.

MUMPS는 다중전면(multifrontal) 알고리즘을 사용하여 fill-in을 최소화하는 재정렬(ordering)을 자동으로 수행한다. 본 문제(변수 수 $n \leq 551$)와 같은 중규모 NLP에서는 직접 솔버가 반복적 솔버(iterative solver)보다 안정적이고 효율적이다.

IPOPT는 $\mathbf{D}$의 관성(inertia)을 검사하여 헤시안의 양정치성을 간접적으로 확인한다. 관성이 올바르지 않으면 정규화 항 $\delta_w \mathbf{I}$를 헤시안에 추가하고 인수분해를 반복한다.

\subsection{CasADi Opti Stack과 자동 미분}

CasADi는 심볼릭 프레임워크를 통해 NLP를 자동으로 구축한다~\cite{Andersson2019}. Opti stack은 다음 과정을 자동화한다.

\begin{enumerate}[leftmargin=*]
  \item \textbf{심볼릭 그래프 구축}: \texttt{opti.variable()}, \texttt{opti.minimize()}, \texttt{opti.subject\_to()}를 통해 결정변수, 목적함수, 제약조건의 심볼릭 표현을 생성한다.
  \item \textbf{자동 미분(Automatic Differentiation, AD)}: 순전파(forward) 및 역전파(reverse) AD를 통해 야코비안과 헤시안을 정확하게 계산한다. 유한 차분 근사와 달리, 기계 정밀도 수준의 도함수를 얻는다.
  \item \textbf{희소성 탐지}: 심볼릭 그래프를 분석하여 야코비안 및 헤시안의 희소 구조를 자동으로 파악한다. Direct collocation 문제에서 야코비안은 대역(banded) 구조를 가지며, CasADi는 이를 활용하여 0이 아닌 원소만 계산한다.
  \item \textbf{NLP 인터페이스}: 구축된 심볼릭 NLP를 IPOPT의 C++ 인터페이스에 전달한다. 솔버가 야코비안/헤시안을 요청할 때마다 CasADi가 AD를 통해 효율적으로 평가한다.
\end{enumerate}

이러한 자동화 덕분에, 사용자는 수학적 정식화에 집중할 수 있으며 도함수를 수동으로 유도할 필요가 없다.

\subsection{Two-Pass 전략의 수학적 근거}

Two-Pass 전략은 허용 오차 연속체(tolerance continuation)의 한 형태이다. 일반적으로 NLP의 수렴은 초기값의 품질과 허용 오차 수준에 민감하다. 느슨한 허용 오차에서 먼저 대략적인 해를 구하고, 이를 초기값으로 사용하여 엄격한 허용 오차에서 다시 풀면, 수렴 가능성이 향상된다~\cite{Betts2010}.

\textbf{Pass~1 (Hermite-Simpson, 탐색 단계)}:

\begin{itemize}[leftmargin=*]
  \item 허용 오차: $\epsilon_1 = 10^{-4}$
  \item 목적: 추력 프로파일의 전역적 구조(피크 위치, 개수) 파악
  \item 균일 격자($M = 30$)로 전체 시간 구간을 커버
\end{itemize}

\textbf{Pass~2 (Multi-Phase LGL, 정밀화 단계)}:

\begin{itemize}[leftmargin=*]
  \item 허용 오차: $\epsilon_2 = 10^{-6}$ ($= \epsilon_1 / 100$)
  \item 목적: Pass~1에서 발견한 추력 프로파일 구조를 기반으로 정밀 최적화
  \item 적응적 격자: 피크 구간에 15개, coasting 구간에 8개 LGL 노드 배치
  \item Warm start: Pass~1 해를 보간하여 초기값으로 사용
\end{itemize}

수학적으로, $\epsilon_1 > \epsilon_2$인 두 허용 오차에서의 해 $\mathbf{z}_1^*$, $\mathbf{z}_2^*$ 사이의 관계는 다음과 같다. $\epsilon_1$-최적해 $\mathbf{z}_1^*$가 $\epsilon_2$-최적해 $\mathbf{z}_2^*$의 수렴 반경 내에 있을 확률이, 임의의 초기값을 사용하는 것보다 높다. 이는 해의 연속 의존성(continuous dependence)에 근거한다.

\subsection{수렴 실패 복구의 수학적 근거}

궤적 최적화 NLP는 비볼록(non-convex)이므로, 초기값에 따라 국소 최적해가 달라지거나 수렴에 실패할 수 있다. \texttt{TwoPassOptimizer}는 다음 세 단계의 복구 전략을 순차적으로 적용한다.

\textbf{(1) 다중 시작점 전략 (Multi-start).}
%
\begin{equation}\label{eq:multi-start}
\nu_0^{(k)} \sim \mathcal{U}(0, 2\pi), \quad \nu_f^{(k)} \sim \mathcal{U}(0, 2\pi), \quad k = 1, 2, 3
\end{equation}

True anomaly $\nu_0$, $\nu_f$를 무작위로 변경하여 최대 3회 재시도한다. 경계조건이 $\nu_0$, $\nu_f$에 비선형적으로 의존하므로, 다른 $\nu$ 값은 NLP의 실행가능 영역의 기하학적 구조를 변경하여 새로운 수렴 경로를 열 수 있다. 원궤도에서 $\nu$는 argument of latitude와 동치이므로, $[0, 2\pi)$에서의 균일 샘플링은 궤도상의 모든 위치를 동등하게 탐색한다.

재현성을 위해 난수 생성기의 시드를 42로 고정한다(\texttt{np.random.default\_rng(42)}).

\textbf{(2) 허용 오차 완화 (Tolerance relaxation).}
%
\begin{equation}\label{eq:tol-relaxation}
\epsilon_{\text{relaxed}} = \epsilon_1 \times \alpha, \quad \alpha = 10
\end{equation}

허용 오차를 $10^{-4}$에서 $10^{-3}$으로 10배 완화한다. 이는 실행가능성 복원(feasibility restoration)의 한 형태로 볼 수 있다. 더 넓은 허용 오차는 KKT 조건의 만족 범위를 확대하여, 엄격한 조건에서는 수렴하지 않았을 해를 찾을 수 있게 한다.

\textbf{(3) 수렴 실패 반환.}

모든 복구 시도가 실패하면, 최종 반복의 해를 \texttt{converged = False}로 표시하여 반환한다. 비용은 $+\infty$로 설정되어, 후속 분석에서 해당 케이스가 유효하지 않음을 나타낸다.


%% ============================================================================
\section{구현 매핑}
%% ============================================================================

\subsection{solver.py: IPOPT 옵션 관리}

\texttt{solver.py}(32행)는 두 개의 함수 \texttt{get\_pass1\_options}과 \texttt{get\_pass2\_options}으로 구성되며, Pass별 IPOPT 옵션 딕셔너리를 반환한다.

\begin{table}[h]
\centering
\begin{tabular}{@{}llll@{}}
\toprule
함수 & 행 & 기본 옵션 원천 & 용도 \\
\midrule
\texttt{get\_pass1\_options(**overrides)} & 6--17행 & \texttt{IPOPT\_OPTIONS\_PASS1} & H-S collocation 솔버 설정 \\
\texttt{get\_pass2\_options(**overrides)} & 20--31행 & \texttt{IPOPT\_OPTIONS\_PASS2} & Multi-Phase LGL 솔버 설정 \\
\bottomrule
\end{tabular}
\end{table}

두 함수 모두 \texttt{**overrides} 키워드 인자를 받아 기본 설정을 덮어쓸 수 있다. 이 설계는 기본 설정과 케이스별 예외 설정을 분리하며, 수렴 실패 복구 시 특정 옵션(예: \texttt{tol})만 선택적으로 변경하는 데 사용된다.

\subsection{IPOPT 옵션 상세}

Pass~1과 Pass~2의 IPOPT 옵션은 \texttt{config.py}에 정의되어 있다.

\begin{table}[h]
\centering
\begin{tabular}{@{}llll@{}}
\toprule
파라미터 & Pass~1 & Pass~2 & 의미 \\
\midrule
\texttt{tol} & $10^{-4}$ & $10^{-6}$ & 전체 KKT 오차 수렴 기준 \\
\texttt{constr\_viol\_tol} & $10^{-4}$ & $10^{-6}$ & 제약조건 위반 허용 오차 \\
\texttt{max\_iter} & 500 & 1000 & 최대 반복 횟수 \\
\texttt{linear\_solver} & MUMPS & MUMPS & 선형 시스템 직접 솔버 \\
\texttt{mu\_strategy} & adaptive & adaptive & 장벽 파라미터 감소 전략 \\
\texttt{warm\_start\_init\_point} & -- & yes & Warm start 활성화 \\
\texttt{warm\_start\_bound\_push} & -- & $10^{-6}$ & Warm start 경계 여유 \\
\texttt{print\_level} & 0 & 0 & 출력 억제 \\
\bottomrule
\end{tabular}
\end{table}

\textbf{수렴 허용 오차의 단계적 감소.} Pass~1에서 $10^{-4}$, Pass~2에서 $10^{-6}$으로 100배 엄격해진다. 이는 2.11절에서 논의한 허용 오차 연속체의 구현이다.

\textbf{Warm start 설정.} Pass~2는 \texttt{warm\_start\_init\_point = "yes"}로 설정되어, IPOPT가 초기 장벽 파라미터와 슬랙 변수를 사용자 제공 초기값으로부터 추정한다. \texttt{warm\_start\_bound\_push = 1e-6}은 초기값이 변수 경계에 매우 가까울 때의 수치적 문제를 방지하기 위한 최소 여유이다. Warm start는 cold start 대비 반복 횟수를 현저히 줄일 수 있다. Pass~1 해를 보간한 초기값이 이미 실행가능 영역 근방에 있으므로, IPOPT는 실행가능성 복원 단계를 건너뛰고 곧바로 최적성 개선에 집중할 수 있다.

\subsection{two\_pass.py: Two-Pass 최적화 파이프라인}

\texttt{TwoPassOptimizer} 클래스(15--114행)는 전체 Two-Pass 파이프라인을 구현한다.

\textbf{클래스 구조:}

\begin{lstlisting}
TwoPassOptimizer
+-- __init__(config: TransferConfig)
+-- solve() -> TrajectoryResult
\end{lstlisting}

\textbf{\texttt{solve()} 메서드의 실행 흐름:}

\begin{table}[h]
\centering
\small
\begin{tabular}{@{}llll@{}}
\toprule
단계 & 행 & 동작 & NLP 솔버 호출 \\
\midrule
초기값 생성 & 40--42행 & \texttt{linear\_interpolation\_guess()} & -- \\
Pass~1 풀이 & 44--47행 & H-S collocation NLP 풀이 & IPOPT (tol=$10^{-4}$) \\
복구: $\nu$ 변경 & 51--63행 & 최대 3회 재시도 & IPOPT (tol=$10^{-4}$) \\
복구: 허용오차 완화 & 65--78행 & tol=$10^{-3}$으로 재시도 & IPOPT (tol=$10^{-3}$) \\
실패 반환 & 80--82행 & \texttt{converged=False} 반환 & -- \\
피크 탐지 & 87--94행 & 추력 크기 프로파일 분석 & -- \\
Phase 구조 결정 & 92--94행 & 피크 위치/폭 기반 phase 분할 & -- \\
보간 (warm start) & 97--99행 & Pass~1 해 $\to$ Pass~2 격자 보간 & -- \\
Pass~2 풀이 & 102--106행 & Multi-Phase LGL NLP 풀이 & IPOPT (tol=$10^{-6}$, warm start) \\
Pass~2 실패 시 대체 & 109--112행 & Pass~1 결과로 대체 & -- \\
\bottomrule
\end{tabular}
\end{table}

\subsection{수렴 실패 복구 흐름}

수렴 실패 복구는 3단계 위계 구조를 가진다. 각 단계는 이전 단계가 실패할 때만 실행된다.

\textbf{단계 1: 기본 풀이 (44--47행).}

선형 보간 초기값(\texttt{linear\_interpolation\_guess})과 기본 설정으로 Pass~1 NLP를 풀이한다. 대부분의 케이스는 이 단계에서 수렴한다.

\textbf{단계 2: True anomaly 랜덤 변경 (49--63행).}

Pass~1이 수렴하지 않으면, 고정 시드(42)의 균일 분포에서 $(\nu_0, \nu_f)$를 샘플링하여 최대 \texttt{MAX\_NU\_RETRIES = 3}회 재시도한다. 각 시도에서 초기값도 재생성한다.

\begin{lstlisting}[language=Python]
rng = np.random.default_rng(42)
for _ in range(MAX_NU_RETRIES):
    nu0_r = rng.uniform(0, 2 * np.pi)
    nuf_r = rng.uniform(0, 2 * np.pi)
    ...
\end{lstlisting}

\textbf{단계 3: 허용 오차 완화 (65--78행).}

$\nu$ 변경으로도 수렴하지 않으면, 허용 오차를 \texttt{TOL\_RELAXATION\_FACTOR = 10}배 완화한다.

\begin{lstlisting}[language=Python]
relaxed_opts = {
    'ipopt.tol': 1e-4 * TOL_RELAXATION_FACTOR,          # 1e-3
    'ipopt.constr_viol_tol': 1e-4 * TOL_RELAXATION_FACTOR,  # 1e-3
}
\end{lstlisting}

\textbf{단계 4: 최종 실패 (80--82행).}

모든 복구 시도가 실패하면 \texttt{converged = False}, \texttt{pass1\_cost = None}으로 결과를 반환한다.

\subsection{Pass~1에서 Pass~2로의 전환}

Pass~1이 수렴하면, 추력 크기 프로파일에서 피크를 탐지하고 phase 구조를 결정한다.

\begin{lstlisting}[language=Python]
u_mag = np.linalg.norm(result1.u, axis=0)
n_peaks, peak_times, peak_widths = detect_peaks(
    result1.t, u_mag, self.config.T)
profile_class = classify_profile(n_peaks)
phases = determine_phase_structure(
    peak_times, peak_widths, self.config.T)
\end{lstlisting}

\texttt{phases}는 딕셔너리의 리스트로, 각 phase의 시작/종료 시각과 노드 수를 정의한다. 피크 구간에는 \texttt{LGL\_NODES\_PEAK = 15}개, coasting 구간에는 \texttt{LGL\_NODES\_COAST = 8}개의 LGL 노드가 배정된다.

Pass~1 해를 Pass~2 격자로 보간하는 \texttt{interpolate\_pass1\_to\_pass2} 함수가 warm start 초기값을 생성한다. 이 보간 과정은 Hermite-Simpson의 균일 격자에서 LGL의 비균일 격자로의 사상(mapping)이다.

\subsection{Pass~2 실패 시 대체 전략}

Pass~2(LGL)가 수렴하지 않으면, Pass~1 결과를 최종 해로 대체한다(109--112행). 이는 ``정밀화에 실패하더라도 대략적인 해를 확보한다''는 실용적 전략이다. Pass~1 비용(\texttt{pass1\_cost})은 두 Pass의 비용 비교를 위해 별도로 기록된다.

\begin{lstlisting}[language=Python]
if not result2.converged:
    result1.pass1_cost = pass1_cost
    return result1
\end{lstlisting}


%% ============================================================================
\section{수치 검증}
%% ============================================================================

\subsection{IPOPT 수렴 판정}

IPOPT는 다음 세 조건을 모두 만족할 때 수렴을 선언한다~\cite{Wachter2006}.
%
\begin{equation}\label{eq:ipopt-convergence}
\max\left(\frac{\|\nabla_\mathbf{z} \mathcal{L}\|_\infty}{s_d},\; \frac{\|\mathbf{c}\|_\infty}{s_c},\; \frac{\mu_j s_j - \mu}{s_c}\right) \leq \epsilon_{\text{tol}}
\end{equation}

여기서 $s_d$, $s_c$는 문제 크기에 의존하는 스케일링 인자이다. \texttt{tol = 1e-4} (Pass~1) 또는 \texttt{tol = 1e-6} (Pass~2)가 $\epsilon_{\text{tol}}$에 해당한다.

또한, 제약조건 위반의 무한대 노름이 별도의 허용 오차 이하여야 한다.
%
\begin{equation}\label{eq:constr-viol}
\|\mathbf{c}\|_\infty \leq \epsilon_{\text{constr}}, \quad \min_j d_j(\mathbf{z}) \geq -\epsilon_{\text{constr}}
\end{equation}

여기서 $\epsilon_{\text{constr}}$은 \texttt{constr\_viol\_tol}에 해당한다.

\subsection{레퍼런스 케이스 수렴 검증}

테스트 코드(\texttt{tests/})에서 다음 레퍼런스 케이스의 수렴을 검증한다.

\begin{table}[h]
\centering
\begin{tabular}{@{}lllll@{}}
\toprule
케이스 & $\Delta a$ [km] & $\Delta i$ [deg] & $T/T_0$ & 검증 항목 \\
\midrule
R1 & $+200$ & 0 & 2.0 & 수렴, $J > 0$, 추력 상한 \\
R2 & $+500$ & 5 & 3.0 & 수렴, 복합 기동 \\
R3 & 0 & 10 & 2.5 & 수렴, 순수 면변환 \\
R4 & $-200$ & 0 & 1.5 & 수렴, 궤도 하강 \\
R5 & $+1000$ & 10 & 4.0 & 수렴, 대규모 전이 \\
\bottomrule
\end{tabular}
\end{table}

각 케이스에서 확인하는 항목은 다음과 같다.

\begin{enumerate}[leftmargin=*]
  \item \textbf{수렴 여부}: \texttt{result.converged == True}
  \item \textbf{비용 유효성}: $0 < J < \infty$
  \item \textbf{추력 상한 만족}: $\|\mathbf{u}_j\| \leq u_{\max} + 10^{-6}$, $\forall j$
  \item \textbf{최소 고도 만족}: $\|\mathbf{r}_j\| - R_E \geq h_{\min} - 1.0$\,km, $\forall j$
  \item \textbf{경계조건 만족}: 출발/도착 궤도요소와의 일치 (제약조건 위반 허용 오차 이내)
\end{enumerate}

\subsection{Two-Pass 수렴 성능}

Two-Pass 전략의 효과를 다음 지표로 평가한다.

\textbf{(1) Pass~2 비용 대 Pass~1 비용 비율.}
%
\begin{equation}\label{eq:cost-ratio}
\rho = \frac{J_2}{J_1}
\end{equation}

$\rho < 1$이면 Pass~2가 더 나은 해를 찾은 것이다. 적응적 격자 배치(피크 구간에 노드 집중)와 더 엄격한 허용 오차 덕분에, 대부분의 케이스에서 $\rho < 1$을 관측한다.

\textbf{(2) 수렴 실패율.} 파라미터 공간 전체에서 최종 수렴 실패(\texttt{converged = False}) 비율을 모니터링한다. 복구 전략이 적용된 후에도 수렴하지 않는 케이스는 파라미터 공간의 경계 영역(매우 짧은 전이 시간, 매우 큰 면변환 등)에 집중되는 경향이 있다.

\subsection{수치 안정성 검증}

\textbf{(1) 야코비안 희소성.} CasADi의 심볼릭 분석에 의해 야코비안의 비영 원소(nonzero) 비율을 확인한다. Hermite-Simpson collocation의 야코비안은 대역 구조를 가지며, 비영 원소 비율이 약 5--10\%로 희소하다. 이 희소성이 MUMPS의 효율적 인수분해를 가능하게 한다.

\textbf{(2) KKT 행렬의 관성.} IPOPT는 매 반복에서 축소 KKT 행렬의 관성(양수/음수/영 고유값의 수)을 검사한다. 올바른 관성은 $(n, m_e, 0)$이다. 관성이 틀리면 정규화를 추가하고 재인수분해한다. 로그에서 ``Number of trial factorizations''가 1이면 관성 보정 없이 진행된 것이다.

\textbf{(3) Warm start 효과.} Pass~2에서 warm start가 활성화되면, IPOPT의 초기 장벽 파라미터 $\mu_0$가 cold start보다 작은 값으로 설정된다. 이는 초기 반복에서 이미 실행가능 영역 근방에 있으므로, 장벽을 빠르게 줄여도 안전하기 때문이다.


%% ============================================================================
\section{참고문헌}
\label{sec:references}
%% ============================================================================

\bibliographystyle{IEEEtran}
\bibliography{references}

\end{document}
